% 目次（PDF 第 12 页）
\BookPage{i}{12}
\BookHeading{目次}{目\quad 次}{}
\BookEntry{7}{iii}{序}{阎崇年}
\BookEntry{9}{v}{捷公祝寿集序言}{冯尔康}
\BookEntry{14}{1}{长松献寿繁花颂}{王耀庭}
\BookEntry{30}{17}{清代中期扬州市镇经济水平和民众生活初探—以刑科题本档案资料为基础}{王跃生}
\BookEntry{50}{37}{清代台湾府与扬州府的关系}{尹章义}
\BookEntry{78}{65}{扬州二十四桥考}{何庄}
\BookEntry{86}{73}{盛清扬州「名医」李炳的医疗生涯及其历史记忆}{余新忠}
\BookEntry{106}{93}{中国社会思想视野中的「难得糊涂」和「郑板桥现象」}{吴建华}
\BookEntry{112}{99}{清乾隆帝南巡与扬州}{岑大利}
\BookEntry{120}{107}{清代扬州的盐务官述略}{杜家骥}
\BookEntry{130}{117}{《条款官话》初探}{赤岭守}
\BookEntry{144}{131}{十七世纪民间天道观研究—基于《广东新语》的研读}{周正庆}
\BookEntry{152}{139}{黄河远上白云间}{林天人}
\BookEntry{174}{161}{清代籍没律例与抄家案件}{柏桦}

% 目次（PDF 第 13 页）
\BookPage{ii}{13}
\BookEntry{190}{177}{翠华南幸扬州写真—盛清君臣眼中的扬州}{庄吉发}
\BookEntry{240}{227}{盛清扬州的城市生活—以《扬州画舫录》为中心}{常建华}
\BookEntry{272}{259}{论清代回疆边外浩罕强权崛起的中国因素}{陈旺城、段昌国}
\BookEntry{282}{269}{康熙小臣养心殿总监造赵昌生平小考}{陈国栋}
\BookEntry{324}{311}{郑板桥对佛教的态度}{陈进传}
\BookEntry{348}{335}{试析雍正乾隆父子对待宗族的差异}{陈龙贵}
\BookEntry{360}{347}{红票：一封康熙皇帝寄给罗马教皇的信}{冯明珠}
\BookEntry{384}{371}{年年红药为谁开—温史话扬州}{冯尔健}
\BookEntry{400}{387}{清代乾隆时期扬州人的引领时尚}{冯尔康}
\BookEntry{434}{421}{记一件康熙朝玻璃器的最高成就}{嵇若昕}
\BookEntry{452}{439}{康雍干时期台湾的乡治}{杨晋平}
\BookEntry{504}{491}{何震年表}{万仕国}
\BookEntry{554}{541}{盛清社会与扬州文化}{阎崇年}
\BookEntry{570}{557}{乾隆及其扬州诗作}{韩小敏}
\BookEntry{590}{577}{乾隆年间的盐商与内帑}{赖惠敏}
\BookEntry{626}{613}{读文渊阁《四库全书》本《明史‧土司传考证》}{戴晋新}
